% ==============================================================================
% PLANTILLA MAESTRA GENERAL - UCNL
% ==============================================================================
% Uso:
% 1. Duplica este archivo y renombra la copia por materia, semana y actividad.
% 2. Actualiza documenttitle, documentsubtitle, documentsubject, coursename,
%    coursecode, universitydepartment y datos de la figura docente.
% 3. Lee la planeacion semanal y NO pegues las instrucciones completas dentro del
%    producto final. Usa la planeacion como lista de cumplimiento.
% 4. Revisa la bibliografia indicada, toma notas, registra citas y redacta con
%    comprension propia.
% 5. Identifica el producto academico solicitado y construyelo dentro del
%    documento siempre que sea posible.
% 6. Entrega en PDF con la nomenclatura indicada en el aula.
%
% Formato institucional recomendado:
% - Fuente tipo Helvetica, equivalente visual a Arial.
% - Interlineado 1.5.
% - Citas autor-anio con natbib: usar \citep{clave} o \citet{clave}.
% - Referencias en bibliografia BibTeX, formato APA 7 en cuerpo y referencias.
% - Caratula, resumen breve, desarrollo, producto solicitado, conclusion y
%   bibliografia.
% - Redaccion academica clara, coherente, cohesionada, formal y sin faltas.
% - La portada puede llevar una marca de agua institucional discreta; ajustar
%   las variables coverwatermark... si se requiere apagarla o cambiar intensidad.
%
% Criterios generales de cumplimiento:
% - Responder exactamente a la actividad solicitada en la planeacion.
% - Identificar conceptos, elementos, criterios o categorias de forma precisa.
% - Analizar, interpretar y tomar postura; no solo copiar fuentes.
% - Organizar la informacion por jerarquia, secuencia y criterios claros.
% - Incluir conclusion personal cuando la actividad lo permita o lo solicite.
% - Usar citas y referencias en APA 7.
% - Evitar copiar las indicaciones de la actividad en el cuerpo final.
%
% Regla para productos visuales:
% - Si la actividad pide tabla, cuadro comparativo, esquema, mapa conceptual,
%   mapa mental, linea de tiempo, diagrama, matriz, lista de verificacion, red
%   conceptual, SQA, organizador grafico o producto similar, realizarlo
%   preferentemente con LaTeX/TikZ para mantener nitidez, evidencia editable y
%   diseno controlado.
% - Si la consigna exige Canva, Miro u otra herramienta externa, insertar una
%   captura nitida y conservar la explicacion dentro del documento.
% - Para tablas extensas, usar pagina limpia, sin encabezados ni pies, landscape,
%   letra pequena y restaurar el estilo al terminar.
%
% Criterios de redaccion personal:
% - No escribir para llenar espacio; escribir para sostener una idea.
% - Convertir informacion en criterio: que significa, como funciona, que problema
%   resuelve, que limite presenta y que consecuencia produce.
% - Estructura mental recomendada: problema -> sistema -> consecuencia -> postura.
% - Usar pensamiento de diagnostico: entrada -> proceso -> salida.
% - Cada parrafo debe tener funcion: presentar, explicar, comparar, valorar o
%   concluir.
% - Usar primera persona con moderacion academica: "Considero que",
%   "Desde esta perspectiva", "La posicion que sostengo es". Evitar "yo creo".
%
% Notas de enfoque personal segun enfoque_personal.txt:
% - La voz debe ser academica, tecnica, directa, estrategica y funcional.
% - El texto no debe sonar como estudiante pasivo que recopila: debe sonar como
%   alguien que interpreta, ordena, cruza datos y busca consecuencia practica.
% - Antes de redactar, identificar la funcion de cada seccion: abrir problema,
%   definir criterio, demostrar con evidencia, comparar, valorar o concluir.
% - Evitar frases genericas como "es importante porque si" o "desde siempre".
%   Sustituirlas por relaciones verificables: causa, efecto, limite, costo,
%   funcion, riesgo, implicacion o mejora.
% - Cada concepto debe pasar de definicion a analisis. Formula util:
%   concepto -> funcion -> riesgo si se aplica mal -> consecuencia -> postura.
% - La postura personal no debe ser opinion suelta. Debe tener tres piezas:
%   posicion, razon y efecto. Ejemplo: "Considero que X, porque Y; por ello Z".
% - Integrar experiencia o criterio propio sin perder formalidad: no narrar la
%   vida personal, sino usar mirada practica para explicar por que el tema
%   importa en educacion, derecho, instituciones, comunicacion o profesion.
%
% Estilo observado en actividades ya realizadas:
% - Abrir con una tesis concreta:
%   "El Derecho no solo ordena conductas: organiza poder, define limites y
%   distribuye consecuencias."
% - Formular el problema en negativo y positivo:
%   "El problema no es solo..., sino..."
% - Convertir el producto visual en argumento:
%   "El cuadro no solo recopila datos; reconstruye la funcion, el problema y el
%   punto que requiere correccion."
% - Cerrar con aprendizaje transferible:
%   "No basta con saber que dice la norma; tambien debe analizarse a quien
%   protege, a quien excluye y que razones permiten defenderla publicamente."
%
% Nivel cognitivo esperado:
% - Recordar: usar solo para definir conceptos basicos. No debe ser el nivel
%   dominante en una actividad ponderable.
% - Comprender: explicar con palabras propias y relacionar con el tema.
% - Aplicar: llevar el concepto a un caso, norma, texto, correo, dilema o hecho.
% - Analizar: separar elementos, detectar causas, efectos, tensiones y limites.
% - Evaluar: sostener una postura con criterios, fuentes y consecuencias.
% - Crear: elaborar producto propio: mapa, cuadro, linea de tiempo, ensayo,
%   propuesta, reformulacion, glosario, portafolio o solucion argumentada.
%
% Regla de nivel minimo:
% - Si la actividad vale puntos y pide producto, no quedarse en recordar o
%   comprender. Debe llegar por lo menos a aplicar y analizar; si pide conclusion
%   o postura, debe llegar a evaluar.
%
% Uso de las 100 tecnicas didacticas:
% - Esta plantilla conserva el manual "Tecnicas didacticas del aprendizaje" que
%   se incorporo en las plantillas de Filosofia del Derecho, Etica y Moral
%   juridica, y Redaccion en contextos virtuales.
% - Cuando la planeacion mencione una tecnica de la coleccion UnADM, identificar:
%   definicion de la tecnica, producto esperado, pasos de construccion, criterios
%   de revision y cierre reflexivo.
% - La tecnica didactica no debe verse como adorno. Debe mostrar seleccion,
%   jerarquia, relacion entre ideas y lectura propia.
% - Siempre que sea viable, construir la tecnica en LaTeX/TikZ: cuadro
%   comparativo, mapa conceptual, esquema, mapa mental, linea de tiempo, matriz,
%   red conceptual, diagrama de flujo, tabla SQA o lista de verificacion.
%
% Checklist antes de entregar:
% [ ] El titulo, materia, semana, docente y fecha estan actualizados.
% [ ] El objetivo coincide con la planeacion.
% [ ] El producto solicitado aparece visible, rotulado y completo.
% [ ] La introduccion plantea un problema o postura, no solo anuncia el tema.
% [ ] El desarrollo explica, conecta y analiza.
% [ ] Hay citas APA autor-anio dentro del texto.
% [ ] La bibliografia final coincide con las fuentes citadas.
% [ ] La conclusion responde que se comprendio y que postura se sostiene.
% [ ] El nivel cognitivo llega al menos a aplicacion y analisis.
% [ ] Si hay conclusion/postura, incluye posicion, razon y consecuencia.
% [ ] La tecnica didactica solicitada esta construida y explicada.
% [ ] No aparecen instrucciones copiadas de la planeacion dentro del producto.
% [ ] Ortografia, acentos, sintaxis y nombres propios revisados.
% [ ] Si se uso IA como apoyo, el texto fue comprendido, revisado y apropiado.
% ==============================================================================

% CREACION DEL DOCUMENTO
\documentclass[
	spanish,
	letterpaper, oneside
]{article}

% ==============================================================================
% INFORMACION DEL DOCUMENTO
% Actualizar en cada actividad.
% Ejemplos:
% \def\documenttitle {Cuadro comparativo de conceptos fundamentales}
% \def\documentsubtitle {Actividad 1 - Ingles I}
% \def\documentsubject {Actividad 1}
% ==============================================================================
\def\documenttitle {Actividad 1 - Ingles I}
\def\documentsubtitle {Actividad 1 - Ingles I}
\def\documentsubject {Actividad 1}

\def\documentauthor {Martin Jonathan de la Cruz}
\def\coursename {Ingles I}
\def\coursecode {ING-I}

\def\universityname {Universidad Ciudadana de Nuevo Leon}
\def\universityfaculty {Programa academico UCNL}
\def\universitydepartment {Ingles I}
\def\universitydepartmentimage {departamentos/logo-ucnl}
\def\universitydepartmentimagecfg {height=1.57cm}
\def\universitylocation {Monterrey, Nuevo Leon}

% INTEGRANTES, PROFESORES Y FECHAS
\def\authortable {
	\begin{tabular}{ll}
		\textbf{Alumno:}
		& \begin{tabular}[t]{l}
			 Martin Jonathan de la Cruz \\
		\end{tabular} \\
		Matricula: & UCNL \\

		\textit{Figura docente:}
		& \begin{tabular}[t]{l}
			Nombre \\ Apellido
		\end{tabular} \\
		& \\
		\multicolumn{2}{l}{Fecha de realizacion: \today} \\
		\multicolumn{2}{l}{\universitylocation}
	\end{tabular}
}

% IMPORTACION DEL TEMPLATE BASE
\input{template}

% FORMATO GENERAL
\usepackage{helvet}
\renewcommand{\familydefault}{\sfdefault}

% TABLAS Y PRODUCTOS VISUALES
\usepackage{array}
\usepackage{longtable}
\usepackage{pdflscape}
\usepackage{tikz}
\usetikzlibrary{
	arrows.meta,
	positioning,
	calc,
	matrix,
	fit,
	shapes.geometric,
	shadows.blur
}

% CITAS EN FORMATO AUTOR-ANIO, COMPATIBLE CON APA 7 EN EL CUERPO DEL TEXTO
% Usar:
% - Cita parentetica: \citep{clave}
% - Cita narrativa: \citet{clave}
% - Varias fuentes: \citep{clave1,clave2}
\setcitestyle{authoryear,open={(},close={)}}

% ==============================================================================
% AYUDAS DE REDACCION
% Quitar o comentar \pendiente cuando el documento final este terminado.
% ==============================================================================
\newcommand{\pendiente}[1]{\textcolor{red}{[PENDIENTE: #1]}}

% ==============================================================================
% MARCA DE AGUA EN PORTADA
% - Se aplica solo a la portada porque usa \AddToShipoutPictureBG*.
% - Para apagarla en alguna actividad: \def\coverwatermarkenabled {false}
% - Usar opacidad baja para que no compita con titulo, datos ni logotipo.
% ==============================================================================
\def\coverwatermarkenabled {true}
\def\coverwatermarkimage {img/departamentos/logo-nuevo-leon.png}
\def\coverwatermarkopacity {0.16}
\def\coverwatermarkwidth {0.72\paperwidth}
\def\coverwatermarkxshift {0cm}
\def\coverwatermarkyshift {-0.35cm}

\newcommand{\insertcoverwatermark}{%
	\ifthenelse{\equal{\coverwatermarkenabled}{true}}{%
		\AddToShipoutPictureBG*{%
			\begin{tikzpicture}[remember picture,overlay]
				\node[
					opacity=\coverwatermarkopacity,
					inner sep=0pt,
					xshift=\coverwatermarkxshift,
					yshift=\coverwatermarkyshift
				] at (current page.center) {%
					\includegraphics[width=\coverwatermarkwidth]{\coverwatermarkimage}%
				};
			\end{tikzpicture}%
		}%
	}{}%
}

\begin{document}

% PORTADA
\insertcoverwatermark
\templatePortrait

% CONFIGURACION DE PAGINA Y ENCABEZADOS
\templatePagecfg
\onehalfspacing

% ==============================================================================
% RESUMEN / ABSTRACT
% Debe ser breve. Formula recomendada:
% Tema + objetivo + tecnica/producto + enfoque personal.
% ==============================================================================
\begin{abstractd}
	\textbf{Objetivo:} Resolver un cuestionario de Ingles I centrado en estructuras gramaticales basicas y reconocer la opcion correcta en cada reactivo.

	\textbf{Actividad:} Cuestionario de opcion multiple sobre present simple, present perfect, reported speech, condicionales, comparativos, modal verbs y estructura basica de la oracion.

	\textbf{Desarrollo:} Los reactivos se organizan como ejercicio de discriminacion gramatical, con opciones cercanas entre si para obligar a identificar la forma correcta y descartar errores frecuentes.

	\textbf{Solucion:} Se presenta una clave de respuestas con criterio breve de resolucion para cada pregunta.
\end{abstractd}

% TABLA DE CONTENIDOS - INDICE
\templateIndex

% CONFIGURACIONES FINALES
\templateFinalcfg

% ======================= INICIO DEL DOCUMENTO =======================

% ==============================================================================
% SECCION 1: INTRODUCCION
% Apertura recomendada:
% - Iniciar con una afirmacion clara, no con una frase generica.
% - Traducir la consigna a un problema academico, juridico, etico, comunicativo
%   o profesional, segun la materia.
% - Explicar por que importa.
% - Cerrar anticipando el enfoque personal.
% ==============================================================================
\section{Actividad}

La actividad consiste en resolver un cuestionario de 30 reactivos sobre contenidos fundamentales de Ingles I. Cada pregunta trabaja una regla puntual de uso, formacion o transformacion gramatical, de modo que el ejercicio funcione como practica dirigida y como verificacion de comprension.

% ==============================================================================
% SECCION 2: DESARROLLO / ANALISIS
% Adaptar los subtitulos segun la materia y la planeacion.
%
% Productos posibles:
% - Cuadro comparativo: criterios, evidencia, diferencia, valoracion.
% - Mapa conceptual: concepto central, conceptos secundarios, conectores claros.
% - Linea de tiempo: fecha/etapa, hecho, aporte, consecuencia.
% - Estudio de caso: situacion, elementos, problema, interpretacion, solucion.
% - Lista de verificacion: criterio observable, indicador, cumplimiento, mejora.
% - Glosario: concepto, definicion, funcion, ejemplo, fuente.
% - Reporte o ensayo: tesis, argumentos, evidencia, postura y cierre.
%
% Regla de parrafo:
% Entrada: concepto o problema.
% Proceso: explicacion, fuente y relacion.
% Salida: consecuencia, criterio o postura.
% ==============================================================================
\section{Desarrollo}

\subsection{Organizacion del cuestionario}

El cuestionario se construye con la tecnica didactica de \textbf{cuestionario}, porque organiza preguntas claras y respuestas verificables para revisar comprension de forma directa \citep{unadm100TecnicasDidacticas2023}. Los reactivos abarcan saludos, tiempos verbales, reported speech, comparativos, modal verbs y estructura basica del ingles.

\subsection{Criterio de resolucion}

Para resolver la actividad se identifica primero el tema gramatical de cada reactivo y despues se descartan las opciones que rompen la estructura correcta. El desarrollo del cuestionario se centra en distinguir auxiliares, cambios de tiempo verbal, uso de terminaciones, orden de palabras y funcion de expresiones temporales.

\subsection{Cuestionario}

El siguiente cuestionario concentra la actividad completa.

{%
\renewcommand{\labelenumii}{\(\circ\) \alph{enumii}.}
\begin{enumerate}
	\item Which is the negative form of Present Simple?
	\begin{enumerate}
		\item I am not working
		\item I do not like pizza
		\item She not plays
		\item He playing not
	\end{enumerate}

	\item What does ``Hello'' mean? (\textquestiondown Que significa ``Hello''?)
	\begin{enumerate}
		\item Hola
		\item Casa
		\item Adios
		\item Gracias
	\end{enumerate}

	\item \textquestiondown Cual es el uso principal del Present Perfect Continuous?
	\begin{enumerate}
		\item Acciones en progreso desde el pasado hasta el presente
		\item Acciones futuras
		\item Acciones hipoteticas
		\item Acciones terminadas en el pasado
	\end{enumerate}

	\item What is the backshift of ``is eating''?
	\begin{enumerate}
		\item Was eating
		\item Eats
		\item Will eat
		\item Had eaten
	\end{enumerate}

	\item Which time expression is the correct reported form of ``today''?
	\begin{enumerate}
		\item The next day
		\item Then
		\item The previous day
		\item That day
	\end{enumerate}

	\item What is incorrect about ``She musts be late''?
	\begin{enumerate}
		\item The subject ``she''
		\item The verb ``be''
		\item The ``-s'' on the modal verb
		\item The use of ``must''
	\end{enumerate}

	\item What is the correct structure of the Present Perfect?
	\begin{enumerate}
		\item Subject + had + past participle
		\item Subject + will + verb
		\item Subject + have/has + past participle
		\item Subject + is + verb-ing
	\end{enumerate}

	\item Complete: ``This car is \underline{\hspace{1.2cm}} than that car.'' (Este carro es \underline{\hspace{1.2cm}} que ese carro.)
	\begin{enumerate}
		\item fastest
		\item more fast
		\item faster
		\item fast
	\end{enumerate}

	\item What is the correct transformation for: ``I am hungry,'' he said?
	\begin{enumerate}
		\item He said that he would be hungry.
		\item He said that he has been hungry.
		\item He said that he was hungry.
		\item He said that he had been hungry.
	\end{enumerate}

	\item Choose the correct first conditional: (Escoge el primer condicional correcto)
	\begin{enumerate}
		\item If I study, I pass.
		\item If I studied, I would pass.
		\item If I study, I will pass.
		\item If I will study, I pass.
	\end{enumerate}

	\item Choose the correct form: ``We \underline{\hspace{0.9cm}} eaten before the meeting started.''
	\begin{enumerate}
		\item had
		\item has
		\item will have
		\item have
	\end{enumerate}

	\item What is the difference between Present Perfect and Past Perfect?
	\begin{enumerate}
		\item Present Perfect = future plans; Past Perfect = past result
		\item Present = routine; Past = memory
		\item Present Perfect connects past to now; Past Perfect = past before past
		\item They are used the same way
	\end{enumerate}

	\item What verb tense is best for reports at A1 level?
	\begin{enumerate}
		\item Present Simple
		\item Past Continuous
		\item Future Perfect
		\item Present Perfect
	\end{enumerate}

	\item Which verb form is used after ``been''?
	\begin{enumerate}
		\item Infinitive
		\item Past
		\item Present
		\item -ing
	\end{enumerate}

	\item Which sentence shows future with ``will''? (\textquestiondown Que oracion usa futuro con ``will''?)
	\begin{enumerate}
		\item I will call you tomorrow.
		\item I was eating.
		\item I am eating.
		\item I eat pizza.
	\end{enumerate}

	\item Choose the correct structure.
	\begin{enumerate}
		\item Smartest than
		\item The most smartest
		\item The more smart
		\item The smartest
	\end{enumerate}

	\item What is the superlative of ``bad''?
	\begin{enumerate}
		\item Baddest
		\item Worst
		\item Worse
		\item More bad
	\end{enumerate}

	\item Which of these is a good essay topic?
	\begin{enumerate}
		\item Class attendance report
		\item School lunch menu
		\item My favorite movie
		\item Annual sales report
	\end{enumerate}

	\item What happens to the verb ``study'' in ``I study English'' when reported?
	\begin{enumerate}
		\item Studied
		\item Studies
		\item Is studying
		\item Had studied
	\end{enumerate}

	\item \textquestiondown Que indica la frase ``I have been working all day''?
	\begin{enumerate}
		\item Accion futura
		\item Accion de ayer
		\item Accion con duracion larga
		\item Accion breve
	\end{enumerate}

	\item What is the correct verb tense transformation for: Present Perfect $\rightarrow$ ?
	\begin{enumerate}
		\item Past Perfect
		\item Simple Past
		\item Present Simple
		\item Future
	\end{enumerate}

	\item Which of these verbs is not commonly used in this tense?
	\begin{enumerate}
		\item Study
		\item Read
		\item Know
		\item Work
	\end{enumerate}

	\item What is the function of modal verbs?
	\begin{enumerate}
		\item To replace main verbs entirely
		\item To express time only
		\item To modify the meaning of the main verb
		\item To act as subjects in a sentence
	\end{enumerate}

	\item What comes first in an essay?
	\begin{enumerate}
		\item Title
		\item Conclusion
		\item Body
		\item Introduction
	\end{enumerate}

	\item Choose the correct form for long adjectives.
	\begin{enumerate}
		\item More difficult
		\item Difficulty
		\item Difficultest
		\item Most difficulty
	\end{enumerate}

	\item Which is the correct question in Present Simple?
	\begin{enumerate}
		\item Do you like music?
		\item Are you like music?
		\item You do like music?
		\item You are liking music?
	\end{enumerate}

	\item Which word often goes with Present Simple?
	\begin{enumerate}
		\item At the moment
		\item Now
		\item Always
		\item Right now
	\end{enumerate}

	\item How do we form Present Simple for ``he, she, it''?
	\begin{enumerate}
		\item Add -ing
		\item Add -ed
		\item Add -s or -es
		\item Use will
	\end{enumerate}

	\item Choose the correct form: ``She \underline{\hspace{0.9cm}} lived here for five years.''
	\begin{enumerate}
		\item did
		\item was
		\item had
		\item has
	\end{enumerate}

	\item What is the basic sentence order in English? (\textquestiondown Cual es el orden basico de la oracion en ingles?)
	\begin{enumerate}
		\item Subject + Verb + Object
		\item Verb + Subject + Object
		\item Object + Subject + Verb
		\item Verb + Object + Subject
	\end{enumerate}
\end{enumerate}
}

\section{Solucion}

La solucion del cuestionario se presenta en una clave de respuestas con criterio breve para cada reactivo.

\subsection{Clave de respuestas}

\begingroup
\small
\setlength{\tabcolsep}{5pt}
\renewcommand{\arraystretch}{1.18}
\begin{longtable}{@{}p{0.08\linewidth}p{0.16\linewidth}p{0.68\linewidth}@{}}
\caption{Clave de respuestas del cuestionario de Ingles I.}\label{tab:clave-cuestionario-ingles}\\
\toprule
\textbf{No.} & \textbf{Correcta} & \textbf{Criterio de respuesta} \\
\midrule
\endfirsthead
\toprule
\textbf{No.} & \textbf{Correcta} & \textbf{Criterio de respuesta} \\
\midrule
\endhead
1 & b & La forma negativa del present simple usa do not con el verbo base. \\
2 & a & Hello significa Hola y corresponde a un saludo basico. \\
3 & a & El present perfect continuous expresa una accion iniciada en el pasado que continua hasta el presente. \\
4 & a & En reported speech, is eating cambia a was eating por backshift verbal. \\
5 & d & Today suele transformarse en that day en estilo indirecto. \\
6 & c & Los modal verbs no agregan -s en tercera persona; por eso musts es incorrecto. \\
7 & c & El present perfect se forma con have o has mas participio pasado. \\
8 & c & El comparativo correcto de fast es faster. \\
9 & c & I am hungry pasa a he was hungry en estilo indirecto. \\
10 & c & El first conditional correcto combina if + present simple con will en la oracion principal. \\
11 & a & Before the meeting started exige had eaten para marcar anterioridad respecto de otro pasado. \\
12 & c & El present perfect conecta pasado y presente; el past perfect expresa pasado anterior a otro pasado. \\
13 & a & En nivel A1 el present simple es la base mas estable para reportes breves y claros. \\
14 & d & Despues de been en esta construccion se usa el verbo en forma -ing. \\
15 & a & I will call you tomorrow expresa futuro con will. \\
16 & d & El superlativo correcto es the smartest. \\
17 & b & El superlativo irregular de bad es worst. \\
18 & c & My favorite movie es un tema apropiado para ensayo breve; las otras opciones remiten mas a reporte o lista. \\
19 & a & En reported speech, study cambia a studied. \\
20 & c & Have been working all day indica duracion prolongada de la accion. \\
21 & a & En la secuencia de backshift, present perfect se transforma en past perfect. \\
22 & c & Know es un verbo estativo y no suele usarse en present perfect continuous. \\
23 & c & Los modal verbs modifican el significado del verbo principal. \\
24 & a & En la presentacion de un ensayo, el titulo aparece primero. \\
25 & a & Los adjetivos largos forman el comparativo con more. \\
26 & a & La pregunta correcta en present simple usa auxiliar: Do you like music? \\
27 & c & Always es un adverbio de frecuencia tipico del present simple. \\
28 & c & Con he, she e it se agrega -s o -es al verbo en present simple. \\
29 & d & She has lived here for five years usa present perfect para una situacion que continua. \\
30 & a & El orden basico en ingles es Subject + Verb + Object. \\
\bottomrule
\end{longtable}
\endgroup

El desarrollo queda resuelto con el cuestionario y su clave. La actividad se concentra en identificar la respuesta correcta y justificarla con una regla breve de uso.

% ==============================================================================
% MANUAL DE TECNICAS DIDACTICAS
% Dejar incluido en la plantilla maestra. En una entrega final, puede comentarse
% esta linea si la docente solo solicita el producto de la semana.
% ==============================================================================
% Referencia opcional desactivada: archivo externo no disponible en este repositorio.

% BIBLIOGRAFIA
% Cambiar el nombre del archivo .bib segun la materia/actividad.
% Ejemplo: \bibliography{filosofia-del-derecho}
\bibliography{ingles-I}

% FIN DEL DOCUMENTO
\end{document}


